\documentclass[11pt]{article}
\usepackage[T1]{fontenc}
\usepackage[utf8]{inputenc}
\usepackage{lmodern}
\usepackage{microtype}
\usepackage[margin=1in,headheight=15pt]{geometry}
\usepackage{amsmath,amssymb,amsthm,mathtools}
% Canonical mathematical notation for active FabiusFunction LaTeX documents.
%
% This file is intentionally a plain TeX input rather than a LaTeX package.
% Every standalone document loads it by an explicit path relative to that
% document's directory; included fragments inherit the definitions from their
% root document.  Keep this file independent of document layout and metadata.
%
% Contract:
%   * one command has one mathematical meaning throughout the active corpus;
%   * names encode semantic roles rather than a document-local choice of letter;
%   * visually similar but differently normalized objects have different print;
%   * every command below is defined normatively in the notation catalogue.
%
% The active corpus excludes docs/archive/.  Historical sources may retain
% their original notation only inside an explicitly labelled quotation.

\makeatletter
\@ifundefined{mathbb}{%
  \PackageError{fabius-notation}{The amssymb package must be loaded first}%
    {Load amsmath and amssymb before inputting fabius-notation.tex.}%
}{}
\@ifundefined{operatorname}{%
  \PackageError{fabius-notation}{The amsmath package must be loaded first}%
    {Load amsmath before inputting fabius-notation.tex.}%
}{}
\makeatother

% Version stamp used by the catalogue and by static audits.
\newcommand{\FabiusNotationVersion}{2026-08-31}

% Number systems.  NaturalNumbers is explicitly the nonnegative convention.
\newcommand{\NaturalNumbers}{\mathbb{N}_{0}}
\newcommand{\PositiveIntegers}{\mathbb{N}_{>0}}
\newcommand{\IntegerNumbers}{\mathbb{Z}}
\newcommand{\RationalNumbers}{\mathbb{Q}}
\newcommand{\RealNumbers}{\mathbb{R}}
\newcommand{\ComplexNumbers}{\mathbb{C}}
\newcommand{\ExtendedRealNumbers}{\overline{\mathbb{R}}}

% The three principal Fabius--Rvachev functions.  Subscripts are deliberate:
% a bare F and a calligraphic F have several unrelated uses in the corpus.
\newcommand{\FabiusBounded}{F_{\mathrm{bd}}}
\newcommand{\FabiusGlobal}{F_{\mathrm{gl}}}
\newcommand{\FabiusQuantile}{Q_{F}}
\newcommand{\FabiusClampedQuantile}{Q_{F,\mathrm{cl}}}
\newcommand{\RvachevUp}{\operatorname{up}}

% Constants, elementary operators, and delimiters.
\newcommand{\EulerE}{\mathrm{e}}
\newcommand{\ImaginaryUnit}{\mathrm{i}}
\newcommand{\Differential}{\mathop{}\!\mathrm{d}}
\newcommand{\AbsoluteValue}[1]{\left\lvert #1\right\rvert}
\newcommand{\Norm}[1]{\left\lVert #1\right\rVert}
\newcommand{\Floor}[1]{\left\lfloor #1\right\rfloor}
\newcommand{\Ceiling}[1]{\left\lceil #1\right\rceil}
\newcommand{\FractionalPart}[1]{\left\{#1\right\}}
\newcommand{\PositivePart}[1]{\left(#1\right)_{+}}
\newcommand{\IndicatorOf}[1]{\mathbf{1}_{#1}}
\newcommand{\SetOf}[1]{\left\{#1\right\}}
\newcommand{\InnerProduct}[1]{\left\langle #1\right\rangle}
\newcommand{\CardinalityOf}[1]{\left\lvert #1\right\rvert}
\newcommand{\DefinitionEquals}{\mathrel{\coloneqq}}
\newcommand{\Sign}{\operatorname{sgn}}
\newcommand{\IdentityOperator}{\operatorname{id}}
\newcommand{\SpanOperator}{\operatorname{span}}
\newcommand{\TraceOperator}{\operatorname{tr}}
\newcommand{\RankOperator}{\operatorname{rank}}
\newcommand{\DiagonalOperator}{\operatorname{diag}}
\newcommand{\ResidueOperator}{\operatorname*{Res}}
\newcommand{\LeastCommonMultiple}{\operatorname{lcm}}
\newcommand{\OrderOperator}{\operatorname{ord}}
\newcommand{\PrincipalLogarithm}{\operatorname{Log}}
\newcommand{\FallingFactorial}[2]{\left(#1\right)^{\underline{#2}}}
\newcommand{\RisingFactorial}[2]{\left(#1\right)^{\overline{#2}}}

% Support, topology, convolution, and metric notation.
\newcommand{\SupportOperator}{\operatorname{supp}}
\newcommand{\SupportOf}[1]{\SupportOperator\!\left(#1\right)}
\newcommand{\TopologicalSupportOf}[1]{\operatorname{tsupp}\!\left(#1\right)}
\newcommand{\Convolution}{\mathbin{*}}
\newcommand{\MetricDistance}{\operatorname{dist}}
\newcommand{\TotalVariationDistance}{d_{\mathrm{TV}}}

% Probability.  Equality and convergence in law are not metric distance.
\newcommand{\Probability}{\mathbb{P}}
\newcommand{\Expectation}{\mathbb{E}}
\newcommand{\Variance}{\operatorname{Var}}
\newcommand{\Covariance}{\operatorname{Cov}}
\newcommand{\LawOperator}{\operatorname{Law}}
\newcommand{\LawOf}[1]{\LawOperator\!\left(#1\right)}
\newcommand{\EqualInLaw}{\mathrel{\overset{\mathrm{law}}{=}}}
\newcommand{\ConvergesInLaw}{\mathrel{\xrightarrow{\mathrm{law}}}}

% Fourier and integral transforms.  The printed labels expose normalization.
% FourierTwoPi uses exp(-2*pi*i*x*xi); FourierAngular uses exp(-i*x*omega).
\newcommand{\FourierTwoPi}[1]{\widehat{#1}_{\,2\pi}}
\newcommand{\FourierAngular}[1]{\widehat{#1}_{\,\mathrm{ang}}}
\newcommand{\LaplaceTransformSymbol}{\mathcal{L}}
\newcommand{\MellinTransformSymbol}{\mathcal{M}}
\newcommand{\LaplaceTransformOf}[1]{\LaplaceTransformSymbol\!\left[#1\right]}
\newcommand{\MellinTransformOf}[1]{\MellinTransformSymbol\!\left[#1\right]}
\newcommand{\MomentGeneratingFunction}[1]{M_{#1}}
\newcommand{\CharacteristicFunction}[1]{\varphi_{#1}}

% The two standard sinc functions must never share one symbol.
\newcommand{\SincRad}{\operatorname{sinc}_{\mathrm{rad}}}
\newcommand{\SincPi}{\operatorname{sinc}_{\pi}}

% Binary arithmetic and Thue--Morse notation.
\newcommand{\BinaryDigitSum}{\operatorname{wt}_{2}}
\newcommand{\ThueMorseSignSymbol}{\varepsilon}
\newcommand{\ThueMorseSign}[1]{\varepsilon_{#1}}
\newcommand{\PadicValuation}[1]{\nu_{#1}}
\newcommand{\TwoAdicValuation}{\nu_{2}}

% q-series notation.  Every base is an explicit argument.
\newcommand{\QPochhammer}[3]{\left(#1;#2\right)_{#3}}
\newcommand{\GaussianBinomial}[3]{%
  \genfrac{[}{]}{0pt}{}{#1}{#2}_{#3}%
}
\newcommand{\QInteger}[2]{\left[#1\right]_{#2}}

% Named special functions.
\newcommand{\Polylogarithm}{\operatorname{Li}}
\newcommand{\LambertW}{W}
\newcommand{\GammaFunction}{\Gamma}
\newcommand{\RiemannZeta}{\zeta}

% Asymptotics.  BigO and LittleO are operator tokens for existing formula
% grammar; prefer the At forms so that the limiting regime is printed.
\newcommand{\BigO}{\mathcal{O}}
\newcommand{\LittleO}{o}
\newcommand{\BigOAt}[2]{\mathcal{O}_{#2}\!\left(#1\right)}
\newcommand{\LittleOAt}[2]{o_{#2}\!\left(#1\right)}

\usepackage{booktabs,longtable,array,tabularx}
\usepackage{enumitem}
\usepackage{xcolor}
\usepackage{fancyhdr}
\usepackage{hyperref}
\usepackage[nameinlink,capitalise,noabbrev]{cleveref}
\usepackage{bookmark}
\usepackage{setspace}

\definecolor{linkblue}{RGB}{28,72,130}
\definecolor{theoremblue}{RGB}{32,72,112}
\hypersetup{colorlinks=true,linkcolor=linkblue,citecolor=linkblue,urlcolor=linkblue,
  pdftitle={q-Series from First Principles},pdfauthor={A proof-driven synthesis},
  pdfsubject={q-series, theta functions, partitions, Bailey pairs, Rogers-Ramanujan identities},
  pdfkeywords={q-series, q-Pochhammer, q-binomial, theta, eta, partitions, Rogers-Ramanujan, Andrews-Gordon}}
\pagestyle{fancy}
\fancyhf{}
\fancyhead[L]{\small q-Series from First Principles}
\fancyhead[R]{\small\nouppercase{\leftmark}}
\fancyfoot[C]{\thepage}
\renewcommand{\headrulewidth}{0.4pt}
\setstretch{1.04}
\setlength{\parindent}{1.4em}
\setlength{\parskip}{0.15em}
\setlist{itemsep=0.25em,topsep=0.45em}
\allowdisplaybreaks
\emergencystretch=3em

\newtheoremstyle{blueplain}{7pt}{7pt}{\itshape}{}{\color{theoremblue}\bfseries}{.}{0.5em}{}
\theoremstyle{blueplain}
\newtheorem{theorem}{Theorem}[section]
\newtheorem{proposition}[theorem]{Proposition}
\newtheorem{lemma}[theorem]{Lemma}
\newtheorem{corollary}[theorem]{Corollary}
\newtheorem{identity}[theorem]{Identity}
\newtheoremstyle{bluedef}{7pt}{7pt}{\normalfont}{}{\color{theoremblue}\bfseries}{.}{0.5em}{}
\theoremstyle{bluedef}
\newtheorem{definition}[theorem]{Definition}
\newtheorem{remark}[theorem]{Remark}
\newtheorem{example}[theorem]{Example}

\newcommand{\qfac}[1]{(q;q)_{#1}}
\newcommand{\qinf}{(q;q)_{\infty}}

\title{\vspace{-1.5em}\textbf{q-Series from First Principles}\\[0.35em]
\large Products, Basic Hypergeometric Sums, Theta Functions, Partitions,\\
Bailey Pairs, and the Rogers--Ramanujan World}
\author{A proof-driven synthesis of the supplied reference collection}
\date{30 August 2026}

\begin{document}
\maketitle

\begin{abstract}
This article develops a self-contained route through the principal themes represented in the supplied collection of MathWorld articles and Wolfram notebooks: $q$-Pochhammer symbols, Gaussian polynomials, basic hypergeometric summation, Jacobi and Ramanujan theta functions, the Dedekind eta function, partition generating functions, Rogers--Ramanujan identities and continued fractions, Bailey pairs, and the Andrews--Gordon identities. The emphasis is not on cataloguing formulas but on exposing a small collection of mechanisms that repeatedly generate them. Every theorem used in the main proof chain is proved in the text. In particular, we give complete proofs of the finite and infinite $q$-binomial theorems, Heine's transformation, the $q$-Gauss and $q$-Chu--Vandermonde sums, the terminating $q$-Pfaff--Saalsch\"utz sum, the Jacobi triple product, Euler's pentagonal theorem and partition recurrence, theta and eta transformation laws, Bailey inversion and Bailey's lemma, both Rogers--Ramanujan identities and their continued fraction, and the full Andrews--Gordon family. A final atlas explains how the additional supplied entries---including Fine, Jackson, Jackson--Slater, Rogers--Selberg, Bailey mod $9$, Rogers mod $14$, Dyson mod $27$, quintuple-product, Schr\"oter, Dougall--Ramanujan, and Borwein material---fit into the same architecture.
\end{abstract}

\tableofcontents
\newpage

\section{Purpose, scope, and the supplied corpus}

The supplied archive is unusually coherent. Its files fall into four interacting layers:
\begin{enumerate}[label=\textbf{\arabic*.}]
\item \emph{Algebraic and hypergeometric foundations}: $q$-Pochhammer symbols, $q$-binomial coefficients, basic hypergeometric functions, $q$-Vandermonde and $q$-Saalsch\"utz summations, Fine's equation, and Jackson-type identities.
\item \emph{Product and modular structure}: the Jacobi triple and quintuple products, Jacobi and Ramanujan theta functions, Schr\"oter's formula, and the Dedekind eta function.
\item \emph{Enumerative and arithmetic applications}: the ordinary and distinct-part partition functions, Euler's pentagonal number theorem, and sums of squares.
\item \emph{Rogers--Ramanujan families}: the two classical identities, their continued fraction, Rogers--Selberg and higher-modulus identities, Bailey mod $9$, Rogers mod $14$, Dyson mod $27$, and the Andrews--Gordon theorem.
\end{enumerate}

A literal transcription of the individual reference pages would obscure their common structure. We instead organize the material around the chain
\[
\begin{gathered}
\boxed{\text{finite }q\text{-algebra}}
\Longrightarrow
\boxed{\text{basic hypergeometric transformations}}
\Longrightarrow
\boxed{\text{theta products}}\\[0.5em]
\boxed{\text{theta products}}
\Longrightarrow
\boxed{\text{partitions and modular products}}
\Longrightarrow
\boxed{\text{Bailey chains and lattices}}.
\end{gathered}
\]
The final section returns to every supplied topic and identifies the relevant proof engine.

\subsection{Proof policy}

The phrase ``complete proof'' is used literally in the main line of the article. Interchanges of infinite sums and products are justified either by absolute convergence or by working coefficientwise in a formal power-series ring. Theorems that are only mentioned for orientation in the final source atlas are explicitly labelled as such and are not used in any proof. Throughout the analytic portions we assume $0<|q|<1$, unless a finite identity makes this unnecessary. All denominator parameters are assumed to avoid their evident poles; because the terminating identities are rational identities in the parameters, the generic statements extend to all nonsingular specializations.

\section{Analytic and formal foundations}

\begin{definition}[$q$-Pochhammer symbols]
For $n\in\NaturalNumbers$ define
\[
(a;q)_0=1,\qquad (a;q)_n=\prod_{j=0}^{n-1}(1-aq^j).
\]
For $|q|<1$ define
\[
(a;q)_\infty=\prod_{j=0}^{\infty}(1-aq^j).
\]
Multiple parameters are abbreviated by
\[
(a_1,\ldots,a_r;q)_n=\prod_{j=1}^r(a_j;q)_n,
\qquad
(a_1,\ldots,a_r;q)_\infty=\prod_{j=1}^r(a_j;q)_\infty.
\]
\end{definition}

The elementary shift laws
\begin{align}
(a;q)_{m+n}&=(a;q)_m(aq^m;q)_n, \label{eq:shift-finite}\\
(a;q)_\infty&=(a;q)_n(aq^n;q)_\infty \label{eq:shift-infinite}
\end{align}
will be used constantly.

\begin{lemma}[Absolute convergence of the basic products]\label{lem:product-convergence}
Let $|q|<1$. On every compact set of the $a$-plane, $(a;q)_\infty$ converges uniformly and hence defines an entire function of $a$. If no factor vanishes, the product is nonzero.
\end{lemma}

\begin{proof}
Fix $|a|\le M$. Then
\[
\sum_{j=0}^{\infty}|aq^j|\le \frac{M}{1-|q|}<\infty.
\]
Choose $J$ so large that $|aq^j|\le \tfrac12$ for $j\ge J$, uniformly on the compact set. The principal logarithm satisfies
\[
|\log(1-w)+w|\le C|w|^2\qquad(|w|\le\tfrac12)
\]
for an absolute constant $C$. Thus $\sum_{j\ge J}\log(1-aq^j)$ converges uniformly by the Weierstrass test, so its exponential and the finitely many initial factors give uniform convergence of the product. If no factor is zero, the limiting exponential is nonzero. Uniform convergence on compact sets gives holomorphy; since the compact set was arbitrary, the function is entire.
\end{proof}

\begin{remark}[Formal versus analytic identities]
An identity in $\ComplexNumbers[[q]]$ is established once the coefficient of every $q^N$ agrees. In products such as $\prod_{m\ge1}(1-q^m)^{-1}$, a fixed coefficient involves only finitely many factors. Consequently, most partition identities below are simultaneously formal and analytic. Analytic language is needed chiefly when parameters other than $q$ vary or when modular transformations are considered.
\end{remark}

\section{Gaussian polynomials and finite $q$-algebra}

\begin{definition}[Gaussian polynomial]
For integers $0\le k\le n$, set
\[
\GaussianBinomial{n}{k}{q}=\frac{(q;q)_n}{(q;q)_k(q;q)_{n-k}},
\]
and set it equal to $0$ for $k<0$ or $k>n$.
\end{definition}

\begin{proposition}[$q$-Pascal recurrences]\label{prop:q-pascal}
For $n\ge1$,
\begin{align}
\GaussianBinomial{n}{k}{q}
 &=\GaussianBinomial{n-1}{k}{q}+q^{n-k}\GaussianBinomial{n-1}{k-1}{q}, \label{eq:qpascal1}\\
 &=q^k\GaussianBinomial{n-1}{k}{q}+\GaussianBinomial{n-1}{k-1}{q}. \label{eq:qpascal2}
\end{align}
In particular, all Gaussian coefficients lie in $\IntegerNumbers_{\ge0}[q]$.
\end{proposition}

\begin{proof}
Put the two terms in \eqref{eq:qpascal1} over the common denominator $(q;q)_k(q;q)_{n-k}$. The numerator is
\[
(q;q)_{n-1}\bigl((1-q^{n-k})+q^{n-k}(1-q^k)\bigr)
=(q;q)_{n-1}(1-q^n)=(q;q)_n.
\]
This proves the first recurrence; the second follows similarly or by the symmetry $\GaussianBinomial{n}{k}{q}=\GaussianBinomial{n}{n-k}{q}$. Starting with the boundary values $1$, either recurrence proves polynomiality and coefficientwise nonnegativity by induction on $n$.
\end{proof}

\begin{theorem}[Finite $q$-binomial theorem]\label{thm:finite-qbinomial}
For every $N\in\NaturalNumbers$,
\begin{align}
\prod_{j=0}^{N-1}(1+zq^j)
 &=\sum_{k=0}^{N}q^{\binom{k}{2}}\GaussianBinomial{N}{k}{q} z^k, \label{eq:finite-qbin-plus}\\
(z;q)_N
 &=\sum_{k=0}^{N}(-1)^kq^{\binom{k}{2}}\GaussianBinomial{N}{k}{q} z^k. \label{eq:finite-qbin-minus}
\end{align}
\end{theorem}

\begin{proof}
The second formula follows from the first by replacing $z$ with $-z$. Let $P_N(z)$ denote the right side of \eqref{eq:finite-qbin-plus}. Using \eqref{eq:qpascal1},
\begin{align*}
P_N(z)
&=\sum_k q^{\binom{k}{2}}\GaussianBinomial{N-1}{k}{q}z^k
 +\sum_k q^{\binom{k}{2}+N-k}\GaussianBinomial{N-1}{k-1}{q}z^k\\
&=P_{N-1}(z)+zq^{N-1}P_{N-1}(z)
=(1+zq^{N-1})P_{N-1}(z).
\end{align*}
Since $P_0(z)=1$, induction gives the product.
\end{proof}

\begin{corollary}[Finite reciprocal expansion]\label{cor:finite-reciprocal}
For $N\ge1$, as a formal power series in $z$,
\[
\frac1{(z;q)_N}=\sum_{m=0}^{\infty}\GaussianBinomial{N+m-1}{m}{q}z^m.
\]
\end{corollary}

\begin{proof}
Write the proposed right side as $R_N(z)$. For $N=1$ it is the geometric series. The recurrence \eqref{eq:qpascal2} gives
\begin{align*}
(1-zq^{N-1})R_N(z)
&=\sum_{m\ge0}\left(\GaussianBinomial{N+m-1}{m}{q}-q^{N-1}\GaussianBinomial{N+m-2}{m-1}{q}\right)z^m\\
&=\sum_{m\ge0}\GaussianBinomial{N+m-2}{m}{q}z^m=R_{N-1}(z).
\end{align*}
Induction yields $(z;q)_NR_N(z)=1$.
\end{proof}

The coefficient of $q^m$ in $\GaussianBinomial{n}{k}{q}$ counts partitions of $m$ whose Ferrers diagram fits inside a $k\times(n-k)$ rectangle. Such diagrams are in bijection with lattice paths from $(0,0)$ to $(n-k,k)$, and the recurrences in \cref{prop:q-pascal} separate the paths by their last step. This gives a second, purely combinatorial proof of positivity.

\section{The infinite $q$-binomial theorem and Heine's transformation}

\begin{theorem}[Infinite $q$-binomial theorem]\label{thm:infinite-qbinomial}
If $|q|<1$ and $|z|<1$, then
\begin{equation}\label{eq:infinite-qbinomial}
\sum_{n=0}^{\infty}\frac{(a;q)_n}{(q;q)_n}z^n
=\frac{(az;q)_\infty}{(z;q)_\infty}.
\end{equation}
The equality also holds formally whenever $a$ is treated as an independent parameter.
\end{theorem}

\begin{proof}
Let $F(z)=\sum_{n\ge0}c_nz^n$, where $c_n=(a;q)_n/(q;q)_n$. The ratio $c_n/c_{n-1}=(1-aq^{n-1})/(1-q^n)$ implies, by comparing coefficients,
\begin{equation}\label{eq:qdiff-F}
(1-z)F(z)=(1-az)F(qz),\qquad F(0)=1.
\end{equation}
The product $G(z)=(az;q)_\infty/(z;q)_\infty$ satisfies the same equation by removing the first factors of numerator and denominator, and $G(0)=1$. The equation recursively determines every Taylor coefficient from the preceding one, so $F=G$ as formal series. Analytically, both sides converge normally on compact subsets of $|z|<1$.
\end{proof}

\begin{corollary}[Euler's two expansions]\label{cor:euler-expansions}
For $|z|<1$,
\begin{align}
\frac1{(z;q)_\infty}&=\sum_{n=0}^{\infty}\frac{z^n}{(q;q)_n}, \label{eq:euler-first}\\
(z;q)_\infty&=\sum_{n=0}^{\infty}\frac{(-1)^nq^{\binom n2}}{(q;q)_n}z^n. \label{eq:euler-second}
\end{align}
The second series is entire in $z$.
\end{corollary}

\begin{proof}
Set $a=0$ in \cref{thm:infinite-qbinomial} for \eqref{eq:euler-first}. For \eqref{eq:euler-second}, let $N\to\infty$ in \eqref{eq:finite-qbin-minus}. For fixed $n$,
\[
\GaussianBinomial{N}{n}{q}=\frac{(q^{N-n+1};q)_n}{(q;q)_n}\longrightarrow\frac1{(q;q)_n}.
\]
The convergence is coefficientwise, and normal convergence on compact $z$-sets follows from the quadratic factor $|q|^{\binom n2}$.
\end{proof}

\begin{definition}[Basic hypergeometric series]
With the modern normalization,
\[
{}_r\phi_s\!\left[\begin{matrix}a_1,\ldots,a_r\\ b_1,\ldots,b_s\end{matrix};q,z\right]
=\sum_{n=0}^{\infty}
\frac{(a_1,\ldots,a_r;q)_n}{(q,b_1,\ldots,b_s;q)_n}
\left((-1)^nq^{\binom n2}\right)^{1+s-r}z^n.
\]
When $r=s+1$, the extra factor has exponent zero.
\end{definition}

\begin{theorem}[Heine's first transformation]\label{thm:heine}
Under absolute-convergence conditions, and thereafter by analytic continuation,
\begin{equation}\label{eq:heine}
{}_2\phi_1\!\left[\begin{matrix}a,b\\ c\end{matrix};q,z\right]
=\frac{(b,az;q)_\infty}{(c,z;q)_\infty}
{}_2\phi_1\!\left[\begin{matrix}c/b,z\\ az\end{matrix};q,b\right].
\end{equation}
\end{theorem}

\begin{proof}
By \eqref{eq:shift-infinite} and the infinite $q$-binomial theorem,
\begin{align*}
\frac{(b;q)_n}{(c;q)_n}
&=\frac{(b;q)_\infty}{(c;q)_\infty}\frac{(cq^n;q)_\infty}{(bq^n;q)_\infty}\\
&=\frac{(b;q)_\infty}{(c;q)_\infty}
  \sum_{m\ge0}\frac{(c/b;q)_m}{(q;q)_m}(bq^n)^m.
\end{align*}
Insert this in the defining series and reverse the absolutely convergent sums:
\begin{align*}
{}_2\phi_1(a,b;c;q,z)
&=\frac{(b;q)_\infty}{(c;q)_\infty}
\sum_{m\ge0}\frac{(c/b;q)_m b^m}{(q;q)_m}
\sum_{n\ge0}\frac{(a;q)_n}{(q;q)_n}(zq^m)^n\\
&=\frac{(b;q)_\infty}{(c;q)_\infty}
\sum_{m\ge0}\frac{(c/b;q)_m b^m}{(q;q)_m}
\frac{(azq^m;q)_\infty}{(zq^m;q)_\infty}.
\end{align*}
Finally,
\[
\frac{(azq^m;q)_\infty}{(zq^m;q)_\infty}
=\frac{(az;q)_\infty}{(z;q)_\infty}\frac{(z;q)_m}{(az;q)_m},
\]
which gives \eqref{eq:heine}.
\end{proof}

\begin{corollary}[$q$-Gauss summation]\label{cor:q-gauss}
If $|c/(ab)|<1$, then
\begin{equation}\label{eq:q-gauss}
{}_2\phi_1\!\left[\begin{matrix}a,b\\c\end{matrix};q,\frac{c}{ab}\right]
=\frac{(c/a,c/b;q)_\infty}{(c,c/(ab);q)_\infty}.
\end{equation}
\end{corollary}

\begin{proof}
Put $z=c/(ab)$ in \eqref{eq:heine}. Then $az=c/b$, so one numerator and denominator parameter cancel. Euler's first expansion gives
\[
\sum_{m\ge0}\frac{(c/(ab);q)_m}{(q;q)_m}b^m
=\frac{(c/a;q)_\infty}{(b;q)_\infty}.
\]
Multiplication by the prefactor yields \eqref{eq:q-gauss}.
\end{proof}

\begin{corollary}[$q$-Chu--Vandermonde]\label{cor:q-chu}
For $N\in\NaturalNumbers$,
\begin{equation}\label{eq:q-chu}
{}_2\phi_1\!\left[\begin{matrix}q^{-N},a\\c\end{matrix};q,\frac{cq^N}{a}\right]
=\frac{(c/a;q)_N}{(c;q)_N}.
\end{equation}
\end{corollary}

\begin{proof}
In \eqref{eq:q-gauss}, let $b\to q^{-N}$. The series terminates at $N$, while
\[
\frac{(c/a,cq^N;q)_\infty}{(c,cq^N/a;q)_\infty}
=\frac{(c/a;q)_N}{(c;q)_N}
\]
by \eqref{eq:shift-infinite}. Both sides are rational functions of the remaining parameters, so the identity extends to every nonsingular specialization.
\end{proof}

\section{Euler telescoping and the terminating $q$-Saalsch\"utz sum}

The finite $q$-binomial and $q$-Chu--Vandermonde sums already account for many entries in the supplied notebooks.  The next balanced ${}_3\phi_2$ summation is the algebraic heart of Bailey's lemma.

\begin{lemma}[Euler's telescoping lemma]\label{lem:euler-telescope}
Let $u_k,v_k$ be quantities for which the displayed denominators are nonzero, and put $w_k=u_k-v_k$.  Then
\begin{equation}\label{eq:euler-telescope}
\sum_{k=0}^{N}\frac{w_k}{w_0}
 \frac{u_0u_1\cdots u_{k-1}}{v_1v_2\cdots v_k}
=\frac1{w_0}\left(
 \frac{u_0u_1\cdots u_N}{v_1v_2\cdots v_N}-v_0
\right),
\end{equation}
where empty products are $1$.
\end{lemma}

\begin{proof}
The $k$th summand is
\[
\frac1{w_0}\left(
\frac{u_0\cdots u_k}{v_1\cdots v_k}
-\frac{u_0\cdots u_{k-1}}{v_1\cdots v_{k-1}}
\right),
\]
with the second expression equal to $v_0/w_0$ when $k=0$.  Summation telescopes to \eqref{eq:euler-telescope}.
\end{proof}

\begin{theorem}[$q$-Pfaff--Saalsch\"utz summation]\label{thm:q-pfaff}
For $N\in\NaturalNumbers$,
\begin{equation}\label{eq:q-pfaff}
{}_3\phi_2\!\left[
\begin{matrix}a,b,q^{-N}\\c,abq^{1-N}/c\end{matrix};q,q\right]
=\frac{(c/a,c/b;q)_N}{(c,c/(ab);q)_N}.
\end{equation}
\end{theorem}

\begin{proof}
Normalize the $k$th summand by the reciprocal of the claimed product and set
\begin{equation}\label{eq:F-pfaff}
F(N,k)=
\frac{(c,c/(ab);q)_N}{(c/a,c/b;q)_N}
\frac{(a,b,q^{-N};q)_k}{(c,abq^{1-N}/c,q;q)_k}q^k.
\end{equation}
We prove $\sum_{k=0}^{N}F(N,k)=1$ by induction on $N$.  Extend $F(N,k)$ by zero for $k>N$.  A direct use of the shift law \eqref{eq:shift-finite} gives
\begin{align}
F(N+1,k)-F(N,k)
&=\mathcal C_N
\frac{w_k}{w_0}
\frac{u_0u_1\cdots u_{k-1}}{v_1v_2\cdots v_k}, \label{eq:pfaff-difference}\\
\mathcal C_N
&=-\frac{c(1-a)(1-b)q^N(c,c/(ab);q)_N}
 {ab(c/a,c/b;q)_{N+1}}, \label{eq:pfaff-constant}
\end{align}
where
\begin{align*}
u_k&=(1-aq^k)(1-bq^k)(1-q^{-N-1+k}),\\
v_k&=(1-cq^{k-1})(1-abq^{-N+k}/c)(1-q^k),\\
w_k&=u_k-v_k.
\end{align*}
For completeness, the only non-obvious algebra in \eqref{eq:pfaff-difference} is the elementary identity
\begin{align*}
& (1-cq^N)(1-abq^{-N+k}/c)(1-q^{-N-1})\frac{cq^N}{ab}
 +(1-cq^N/a)(1-cq^N/b)(1-q^{-N+k-1})\\
&\quad=(1-cq^N)(1-c/(abq))(1-q^k)
 +\frac{cq^N}{ab}(1-a)(1-b)(1-q^{-N+k-1}),
\end{align*}
and multiplication by $abq^{-N+k}/c$ converts the right side into $w_k$.  Thus every factor in \eqref{eq:pfaff-difference} has been accounted for.

Now $u_{N+1}=0$ and $v_0=0$.  Applying \cref{lem:euler-telescope} with upper limit $N+1$ therefore gives
\[
\sum_{k=0}^{N+1}\bigl(F(N+1,k)-F(N,k)\bigr)=0.
\]
Hence the normalized sum is independent of $N$.  At $N=0$ it has the single value $F(0,0)=1$, proving \eqref{eq:q-pfaff}.
\end{proof}

\begin{remark}
The proof is a hand-executed creative-telescoping argument.  Its importance for this article is structural: the same balanced sum evaluates the inner finite sum in Bailey's lemma, so the transition from a single hypergeometric identity to infinite families is entirely explicit.
\end{remark}

\section{Partitions, Durfee squares, and Euler products}

\begin{definition}[Partitions]
A partition of $n$ is a weakly decreasing finite sequence $\lambda_1\ge\cdots\ge\lambda_r\ge1$ with sum $n$.  Write $p(n)$ for the number of partitions of $n$, with $p(0)=1$.
\end{definition}

\begin{theorem}[Euler's partition product]\label{thm:partition-product}
As a formal power series, and analytically for $|q|<1$,
\begin{equation}\label{eq:partition-product}
\sum_{n=0}^{\infty}p(n)q^n=\frac1{(q;q)_\infty}.
\end{equation}
\end{theorem}

\begin{proof}
A part of size $m$ may occur $0,1,2,\ldots$ times, contributing the geometric series
\(
1+q^m+q^{2m}+\cdots=(1-q^m)^{-1}.
\)
Multiplying over $m\ge1$ and selecting one term from each factor records every partition exactly once.  For a fixed coefficient $q^n$, factors with $m>n$ contribute only $1$, so the formal product is well defined.
\end{proof}

\begin{proposition}[Partitions in a rectangle]\label{prop:rectangle-partitions}
The generating function for partitions with at most $r$ parts, equivalently with largest part at most $r$, is $1/(q;q)_r$.
\end{proposition}

\begin{proof}
Parts of sizes $1,\ldots,r$ may be repeated freely, giving $1/(q;q)_r$.  Conjugating Ferrers diagrams exchanges ``largest part at most $r$'' with ``at most $r$ parts.''
\end{proof}

\begin{theorem}[Durfee-square decomposition]\label{thm:durfee}
For $|q|<1$,
\begin{equation}\label{eq:durfee}
\frac1{(q;q)_\infty}=\sum_{r=0}^{\infty}\frac{q^{r^2}}{(q;q)_r^2}.
\end{equation}
\end{theorem}

\begin{proof}
Every Ferrers diagram has a unique largest square in its upper-left corner, the Durfee square, say of side $r$.  Removing it leaves a diagram to its right with at most $r$ rows and a diagram below it with largest part at most $r$.  By \cref{prop:rectangle-partitions}, each residual diagram has generating function $1/(q;q)_r$, while the square contributes $q^{r^2}$.  The decomposition is bijective and summing over $r$ gives \eqref{eq:durfee}.
\end{proof}

\begin{theorem}[Partitions into distinct parts]\label{thm:distinct-product}
If $d(n)$ counts partitions of $n$ into distinct parts, then
\begin{equation}\label{eq:distinct-product}
\sum_{n\ge0}d(n)q^n=(-q;q)_\infty.
\end{equation}
If $o(n)$ counts partitions into odd parts, then $d(n)=o(n)$ for every $n$.
\end{theorem}

\begin{proof}
A part $m$ in a distinct-part partition is either absent or present once, giving the factor $1+q^m$ and hence \eqref{eq:distinct-product}.  Moreover,
\[
(-q;q)_\infty
=\prod_{m\ge1}\frac{1-q^{2m}}{1-q^m}
=\frac1{\prod_{m\ge1}(1-q^{2m-1})},
\]
which is exactly the generating function for partitions into odd parts.  Equality of coefficients proves the assertion.
\end{proof}

\section{The Jacobi triple product and its consequences}

\begin{theorem}[Jacobi triple product]\label{thm:jtp}
For $|q|<1$ and $z\ne0$,
\begin{equation}\label{eq:jtp}
\sum_{n\in\IntegerNumbers}(-1)^nq^{\binom n2}z^n
=(q,z,q/z;q)_\infty.
\end{equation}
The series converges normally on compact subsets of $\ComplexNumbers^\times$.
\end{theorem}

\begin{proof}
First consider the annulus $|q|<|z|<1$ and put
\[
F(z)=(z;q)_\infty(q/z;q)_\infty.
\]
It has an absolutely convergent Laurent expansion $F(z)=\sum_{n\in\IntegerNumbers}c_nz^n$.  Removing first factors gives
\[
F(qz)=(qz;q)_\infty(1/z;q)_\infty
=-z^{-1}(z;q)_\infty(q/z;q)_\infty=-z^{-1}F(z).
\]
Comparing coefficients yields $c_{n+1}=-q^nc_n$, and therefore
\begin{equation}\label{eq:jtp-coeff}
c_n=(-1)^nq^{\binom n2}c_0\qquad(n\in\IntegerNumbers).
\end{equation}
It remains to determine $c_0$.  Expanding both factors by Euler's second formula,
\begin{align*}
(z;q)_\infty
 &=\sum_{r\ge0}\frac{(-1)^rq^{\binom r2}}{(q;q)_r}z^r,\\
(q/z;q)_\infty
 &=\sum_{s\ge0}\frac{(-1)^sq^{s(s+1)/2}}{(q;q)_s}z^{-s}.
\end{align*}
The constant term comes from $r=s$, so the Durfee identity gives
\[
c_0=\sum_{r\ge0}\frac{q^{r^2}}{(q;q)_r^2}=\frac1{(q;q)_\infty}.
\]
Substitution in \eqref{eq:jtp-coeff} proves \eqref{eq:jtp} on the annulus.  Both sides are holomorphic on $\ComplexNumbers^\times$ and the series is normally convergent there, so analytic continuation proves the identity for all $z\ne0$.
\end{proof}

\begin{corollary}[Ramanujan's general theta function]\label{cor:ramanujan-theta}
For $|ab|<1$, define
\[
f(a,b)=\sum_{n\in\IntegerNumbers}a^{n(n+1)/2}b^{n(n-1)/2}.
\]
Then
\begin{equation}\label{eq:ramanujan-theta}
f(a,b)=(-a,-b,ab;ab)_\infty.
\end{equation}
\end{corollary}

\begin{proof}
In \eqref{eq:jtp}, replace the base $q$ by $ab$ and put $z=-a$.  The summand becomes
\[
(-1)^n(ab)^{n(n-1)/2}(-a)^n
=a^{n(n+1)/2}b^{n(n-1)/2},
\]
and the product becomes $(-a,-b,ab;ab)_\infty$.
\end{proof}

\begin{corollary}[Euler's pentagonal number theorem]\label{cor:pentagonal}
For $|q|<1$,
\begin{equation}\label{eq:pentagonal}
(q;q)_\infty
=\sum_{n\in\IntegerNumbers}(-1)^nq^{n(3n-1)/2}
=1+\sum_{n=1}^{\infty}(-1)^n
 \left(q^{n(3n-1)/2}+q^{n(3n+1)/2}\right).
\end{equation}
\end{corollary}

\begin{proof}
Apply \eqref{eq:jtp} with base $q^3$ and $z=q$.  The product is
\[
(q^3,q,q^2;q^3)_\infty=(q;q)_\infty,
\]
and the summand is $(-1)^nq^{3n(n-1)/2+n}=(-1)^nq^{n(3n-1)/2}$.  Pairing the terms indexed by $n$ and $-n$ gives the second form.
\end{proof}

\begin{corollary}[Euler's partition recurrence]\label{cor:partition-recurrence}
For $n\ge1$,
\begin{equation}\label{eq:partition-recurrence}
p(n)=\sum_{k=1}^{\infty}(-1)^{k-1}
\left(
 p\!\left(n-\frac{k(3k-1)}2\right)
+p\!\left(n-\frac{k(3k+1)}2\right)
\right),
\end{equation}
where $p(m)=0$ for $m<0$; hence the sum is finite.
\end{corollary}

\begin{proof}
Multiply \eqref{eq:partition-product} by \eqref{eq:pentagonal}; the product is $1$.  Equating the coefficient of $q^n$ for $n>0$ and moving the nonconstant terms to the other side gives \eqref{eq:partition-recurrence}.
\end{proof}

\begin{remark}[Why product identities proliferate]
The triple product converts a bilateral theta series into three arithmetic progressions of Euler factors.  Selecting different bases and values of $z$, combining several instances, or separating residue classes produces the quintuple product, Schr\"oter-type addition formulas, eta quotients, and the product sides of Rogers--Ramanujan-type identities.  The supplied product notebooks are therefore variations on a single analytic mechanism rather than isolated miracles.
\end{remark}

\section{Jacobi theta functions and the Dedekind eta function}

Let $\tau$ lie in the upper half-plane $\mathbb H$ and put $Q=e^{\pi i\tau}$.  We use the theta constants
\begin{align*}
\vartheta_2(\tau)&=\sum_{n\in\IntegerNumbers}Q^{(n+1/2)^2},\\
\vartheta_3(\tau)&=\sum_{n\in\IntegerNumbers}Q^{n^2},\\
\vartheta_4(\tau)&=\sum_{n\in\IntegerNumbers}(-1)^nQ^{n^2},
\end{align*}
and the Dedekind eta function
\begin{equation}\label{eq:eta-def}
\eta(\tau)=e^{\pi i\tau/12}\prod_{n=1}^{\infty}(1-e^{2\pi in\tau})
=Q^{1/12}(Q^2;Q^2)_\infty.
\end{equation}

\begin{theorem}[Theta-product formulas]\label{thm:theta-products}
For $\tau\in\mathbb H$,
\begin{align}
\vartheta_2(\tau)
 &=2Q^{1/4}(Q^2;Q^2)_\infty(-Q^2;Q^2)_\infty^2, \label{eq:theta2-product}\\
\vartheta_3(\tau)
 &=(Q^2;Q^2)_\infty(-Q;Q^2)_\infty^2, \label{eq:theta3-product}\\
\vartheta_4(\tau)
 &=(Q^2;Q^2)_\infty(Q;Q^2)_\infty^2. \label{eq:theta4-product}
\end{align}
Consequently,
\begin{equation}\label{eq:theta-eta-cubic}
\vartheta_2(\tau)\vartheta_3(\tau)\vartheta_4(\tau)=2\eta(\tau)^3.
\end{equation}
\end{theorem}

\begin{proof}
Apply the Jacobi triple product with base $Q^2$.  Taking $z=-Q^2$ gives
\[
\sum_{n\in\IntegerNumbers}Q^{n(n+1)}=(Q^2,-Q^2,-1;Q^2)_\infty
=2(Q^2;Q^2)_\infty(-Q^2;Q^2)_\infty^2.
\]
Multiplication by $Q^{1/4}$ proves \eqref{eq:theta2-product}.  Taking $z=-Q$ and $z=Q$ gives \eqref{eq:theta3-product} and \eqref{eq:theta4-product}.

For the cubic identity, multiply the three products.  The elementary factorizations
\begin{align*}
(Q;Q^2)_\infty(-Q;Q^2)_\infty&=(Q^2;Q^4)_\infty,\\
(-Q^2;Q^2)_\infty&=\frac{(Q^4;Q^4)_\infty}{(Q^2;Q^2)_\infty},\\
(Q^2;Q^2)_\infty&=(Q^2;Q^4)_\infty(Q^4;Q^4)_\infty
\end{align*}
show that all factors apart from $2Q^{1/4}(Q^2;Q^2)_\infty^3$ cancel.  This remaining expression is $2\eta(\tau)^3$ by \eqref{eq:eta-def}.
\end{proof}

\subsection{Poisson summation from first principles}

For an integrable function $f$ on $\RealNumbers$, write
\[
\widehat f(\xi)=\int_{-\infty}^{\infty}f(x)e^{-2\pi i x\xi}\,dx.
\]

\begin{theorem}[Poisson summation for Schwartz functions]\label{thm:poisson}
If $f$ is a Schwartz function, then
\begin{equation}\label{eq:poisson}
\sum_{n\in\IntegerNumbers}f(n)=\sum_{m\in\IntegerNumbers}\widehat f(m),
\end{equation}
and both series converge absolutely.
\end{theorem}

\begin{proof}
The periodization $P(x)=\sum_{n\in\IntegerNumbers}f(x+n)$ and all its derivatives converge uniformly on $[0,1]$, because $f$ is rapidly decreasing.  Thus $P$ is a smooth $1$-periodic function.  Its $m$th Fourier coefficient is
\begin{align*}
\int_0^1P(x)e^{-2\pi imx}\,dx
&=\sum_{n\in\IntegerNumbers}\int_0^1f(x+n)e^{-2\pi imx}\,dx\\
&=\int_{-\infty}^{\infty}f(u)e^{-2\pi imu}\,du
=\widehat f(m).
\end{align*}
Repeated integration by parts makes $\widehat f(m)$ decrease faster than every power of $|m|$, so the Fourier series is absolutely and uniformly convergent.  Evaluating it at $x=0$ gives \eqref{eq:poisson}.
\end{proof}

\begin{lemma}[Fourier transform of a Gaussian]\label{lem:gaussian-ft}
If $\Re t>0$, then, with the branch of $t^{1/2}$ positive for $t>0$,
\begin{equation}\label{eq:gaussian-ft}
\int_{-\infty}^{\infty}e^{-\pi tx^2}e^{-2\pi ix\xi}\,dx
=t^{-1/2}e^{-\pi\xi^2/t}.
\end{equation}
\end{lemma}

\begin{proof}
For real $t>0$, let $I(\xi)$ denote the integral.  Differentiating under the integral and integrating by parts gives
\[
I'(\xi)=-2\pi i\int xe^{-\pi tx^2}e^{-2\pi ix\xi}\,dx
=-\frac{2\pi\xi}{t}I(\xi).
\]
Hence $I(\xi)=I(0)e^{-\pi\xi^2/t}$.  Squaring the elementary Gaussian integral gives $I(0)=t^{-1/2}$.  Both sides of \eqref{eq:gaussian-ft} are holomorphic for $\Re t>0$, so the identity theorem extends the formula to that half-plane.
\end{proof}

\begin{theorem}[Theta transformation laws]\label{thm:theta-transform}
On $\mathbb H$, with the holomorphic square root that is positive at $\tau=i$,
\begin{align}
\vartheta_3(-1/\tau)&=(-i\tau)^{1/2}\vartheta_3(\tau), \label{eq:theta3-S}\\
\vartheta_2(-1/\tau)&=(-i\tau)^{1/2}\vartheta_4(\tau), \label{eq:theta2-S}\\
\vartheta_4(-1/\tau)&=(-i\tau)^{1/2}\vartheta_2(\tau). \label{eq:theta4-S}
\end{align}
\end{theorem}

\begin{proof}
Take $t=-i\tau$, whose real part is positive.  Applying Poisson summation and \cref{lem:gaussian-ft} to $f(x)=e^{-\pi tx^2}$ gives
\[
\sum_{n\in\IntegerNumbers}e^{\pi i\tau n^2}
=(-i\tau)^{-1/2}\sum_{m\in\IntegerNumbers}e^{-\pi im^2/\tau}.
\]
Rearrangement proves \eqref{eq:theta3-S}.

The shifted Poisson formula
\begin{equation}\label{eq:shifted-poisson}
\sum_{n\in\IntegerNumbers}f(n+a)=\sum_{m\in\IntegerNumbers}e^{2\pi ima}\widehat f(m)
\end{equation}
follows by applying the proof of \cref{thm:poisson} to the periodization of $f(x+a)$.  Taking $a=1/2$ yields
\[
\vartheta_2(\tau)=(-i\tau)^{-1/2}\vartheta_4(-1/\tau),
\]
which is equivalent to \eqref{eq:theta4-S}.  Applying ordinary Poisson summation to $e^{2\pi ix(1/2)}f(x)$ gives
\[
\vartheta_4(\tau)=(-i\tau)^{-1/2}\vartheta_2(-1/\tau),
\]
which is \eqref{eq:theta2-S}.
\end{proof}

\begin{theorem}[Dedekind eta transformations]\label{thm:eta-transform}
For $\tau\in\mathbb H$,
\begin{align}
\eta(\tau+1)&=e^{\pi i/12}\eta(\tau), \label{eq:eta-T}\\
\eta(-1/\tau)&=(-i\tau)^{1/2}\eta(\tau). \label{eq:eta-S}
\end{align}
\end{theorem}

\begin{proof}
Equation \eqref{eq:eta-T} follows directly from the product definition: $e^{2\pi in(\tau+1)}=e^{2\pi in\tau}$, while the initial exponential acquires $e^{\pi i/12}$.

For inversion, combine \eqref{eq:theta-eta-cubic} with \cref{thm:theta-transform}:
\[
2\eta(-1/\tau)^3
=(-i\tau)^{3/2}\vartheta_2(\tau)\vartheta_3(\tau)\vartheta_4(\tau)
=2(-i\tau)^{3/2}\eta(\tau)^3.
\]
Thus
\[
R(\tau)=\frac{\eta(-1/\tau)}{(-i\tau)^{1/2}\eta(\tau)}
\]
is holomorphic on the connected domain $\mathbb H$ and satisfies $R(\tau)^3=1$.  Hence $R$ is constant.  At $\tau=i$, inversion fixes $i$ and the chosen square root is $1$, so $R(i)=1$.  This proves \eqref{eq:eta-S}.
\end{proof}

\subsection{Theta powers and sums of squares}

Let $r_d(n)$ denote the number of ordered $d$-tuples $(x_1,\ldots,x_d)\in\IntegerNumbers^d$ with $x_1^2+\cdots+x_d^2=n$.

\begin{proposition}[Theta generating function for sums of squares]\label{prop:squares-theta}
For $|q|<1$,
\begin{equation}\label{eq:squares-theta}
\left(\sum_{m\in\IntegerNumbers}q^{m^2}\right)^d
=\sum_{n=0}^{\infty}r_d(n)q^n.
\end{equation}
\end{proposition}

\begin{proof}
Expand the $d$-fold Cauchy product.  Choosing the term $q^{x_j^2}$ from the $j$th factor produces $q^{x_1^2+\cdots+x_d^2}$, and every ordered representation occurs exactly once.  Absolute convergence justifies the expansion analytically; coefficientwise finiteness proves it formally.
\end{proof}

\begin{remark}
The classical explicit formulas for $r_2(n)$, $r_4(n)$, $r_6(n)$, and $r_8(n)$ arise by expressing powers of theta constants as modular forms or Lambert series.  The supplied ``Sum of Squares Function'' material belongs precisely at this theta-to-arithmetic interface.  Formula \eqref{eq:squares-theta} is the universal generating statement; divisor-sum evaluations require additional modular-form dimension arguments and are not needed later.
\end{remark}

\section{Bailey pairs: inversion, lemma, chain, and lattice shift}

\begin{definition}[Bailey pair]
Two sequences $(\alpha_n)_{n\ge0}$ and $(\beta_n)_{n\ge0}$ form a \emph{Bailey pair relative to $a$} if
\begin{equation}\label{eq:bailey-pair}
\beta_n=\sum_{r=0}^{n}
\frac{\alpha_r}{(q;q)_{n-r}(aq;q)_{n+r}}
\qquad(n\ge0).
\end{equation}
\end{definition}

\subsection{Inverting the Bailey transform}

\begin{lemma}[$q$-difference annihilation]\label{lem:q-annihilation}
If $P(x)$ is a polynomial of degree less than $m$, then
\begin{equation}\label{eq:q-annihilation}
\sum_{s=0}^{m}(-1)^{m-s}q^{\binom{m-s}{2}}
\GaussianBinomial{m}{s}{q} P(q^s)=0.
\end{equation}
\end{lemma}

\begin{proof}
By linearity it suffices to take $P(x)=x^h$ with $0\le h<m$.  Put $t=m-s$ and use the finite $q$-binomial theorem:
\begin{align*}
\sum_{s=0}^{m}(-1)^{m-s}q^{\binom{m-s}{2}}\GaussianBinomial{m}{s}{q} q^{hs}
&=q^{hm}\sum_{t=0}^{m}(-1)^tq^{\binom t2}\GaussianBinomial{m}{t}{q} q^{-ht}\\
&=q^{hm}(q^{-h};q)_m=0,
\end{align*}
because the factor with index $h$ is $1-q^0$.
\end{proof}

\begin{theorem}[Bailey inversion]\label{thm:bailey-inversion}
Equation \eqref{eq:bailey-pair} holds if and only if
\begin{equation}\label{eq:bailey-inversion}
\alpha_n=\frac{1-aq^{2n}}{1-a}
\sum_{j=0}^{n}
\frac{(a;q)_{n+j}}{(q;q)_{n-j}}
(-1)^{n-j}q^{\binom{n-j}{2}}\beta_j.
\end{equation}
The value at $a=1$ is interpreted by its removable limit.
\end{theorem}

\begin{proof}
It is enough to verify that the two triangular matrices are inverse.  Substitute \eqref{eq:bailey-pair} into the right side of \eqref{eq:bailey-inversion}.  The coefficient of $\alpha_r$ is
\begin{equation}\label{eq:inverse-coeff}
C_{n,r}=\frac{1-aq^{2n}}{1-a}
\sum_{j=r}^{n}
\frac{(a;q)_{n+j}(-1)^{n-j}q^{\binom{n-j}{2}}}
{(q;q)_{n-j}(q;q)_{j-r}(aq;q)_{j+r}}.
\end{equation}
If $n=r$, the only term is
\[
\frac{1-aq^{2n}}{1-a}\frac{(a;q)_{2n}}{(aq;q)_{2n}}=1.
\]
Now let $m=n-r>0$ and write $j=r+s$.  The quotient of shifted factorials in \eqref{eq:inverse-coeff} is
\[
\frac{(a;q)_{n+j}}{(aq;q)_{j+r}}
=(1-a)(aq^{2r+s+1};q)_{m-1}.
\]
Therefore
\[
C_{n,r}=\frac{1-aq^{2n}}{(q;q)_m}
\sum_{s=0}^{m}(-1)^{m-s}q^{\binom{m-s}{2}}
\GaussianBinomial{m}{s}{q} P(q^s),
\]
where
\[
P(x)=\prod_{h=0}^{m-2}(1-aq^{2r+1+h}x)
\]
has degree $m-1$.  By \cref{lem:q-annihilation}, $C_{n,r}=0$.  Thus $C_{n,r}=\delta_{n,r}$, proving the inversion.
\end{proof}

\begin{corollary}[The unit Bailey pair]\label{cor:unit-pair}
The sequences
\begin{align}
\beta_n^{\mathrm U}&=\delta_{n,0}, \label{eq:unit-beta}\\
\alpha_n^{\mathrm U}(a)
&=(-1)^nq^{\binom n2}
\frac{1-aq^{2n}}{1-a}\frac{(a;q)_n}{(q;q)_n} \label{eq:unit-alpha}
\end{align}
form a Bailey pair relative to $a$.
\end{corollary}

\begin{proof}
Apply Bailey inversion to $\beta_j=\delta_{j,0}$.  Only the term $j=0$ remains, giving exactly \eqref{eq:unit-alpha}.
\end{proof}

\subsection{Bailey's lemma}

\begin{theorem}[Bailey's lemma]\label{thm:bailey-lemma}
Suppose $(\alpha,\beta)$ is a Bailey pair relative to $a$.  Define
\begin{align}
\alpha_n'&=
\frac{(\rho_1,\rho_2;q)_n}{(aq/\rho_1,aq/\rho_2;q)_n}
\left(\frac{aq}{\rho_1\rho_2}\right)^n\alpha_n,
\label{eq:bailey-alpha-prime}\\
\beta_n'&=\sum_{j=0}^{n}
\frac{(\rho_1,\rho_2;q)_j(aq/(\rho_1\rho_2);q)_{n-j}}
{(q;q)_{n-j}(aq/\rho_1,aq/\rho_2;q)_n}
\left(\frac{aq}{\rho_1\rho_2}\right)^j\beta_j.
\label{eq:bailey-beta-prime}
\end{align}
Then $(\alpha',\beta')$ is also a Bailey pair relative to $a$.
\end{theorem}

\begin{proof}
Put $A=aq$ and $x=A/(\rho_1\rho_2)$.  Substitute the defining expression for $\beta_j$ into \eqref{eq:bailey-beta-prime} and reverse the finite sums.  The coefficient of $\alpha_r$ is
\begin{equation}\label{eq:bailey-inner-before}
\sum_{j=r}^{n}
\frac{(\rho_1,\rho_2;q)_j(x;q)_{n-j}x^j}
{(q;q)_{n-j}(A/\rho_1,A/\rho_2;q)_n
 (q;q)_{j-r}(A;q)_{j+r}}.
\end{equation}
Set $m=n-r$ and $k=j-r$.  We use the elementary reversal identity
\begin{equation}\label{eq:reverse-ratio}
\frac{(x;q)_{m-k}}{(q;q)_{m-k}}
=\frac{(x;q)_m}{(q;q)_m}
\frac{(q^{-m};q)_k}{(q^{1-m}/x;q)_k}
\left(\frac qx\right)^k,
\end{equation}
which follows by writing the last $k$ factors of each finite product in reverse order.  After factoring out the terms independent of $k$, the remaining sum in \eqref{eq:bailey-inner-before} is
\[
{}_3\phi_2\!\left[
\begin{matrix}\rho_1q^r,\rho_2q^r,q^{-m}\\
Aq^{2r},q^{1-m}/x
\end{matrix};q,q\right].
\]
This is balanced, since
\[
\frac{(\rho_1q^r)(\rho_2q^r)q^{1-m}}{Aq^{2r}}
=\frac{q^{1-m}}x.
\]
The $q$-Pfaff--Saalsch\"utz sum therefore evaluates it as
\[
\frac{(Aq^r/\rho_1,Aq^r/\rho_2;q)_m}
{(Aq^{2r},x;q)_m}.
\]
Canceling $(x;q)_m$ and recombining finite products gives the coefficient
\[
\frac{(\rho_1,\rho_2;q)_r x^r}
{(A/\rho_1,A/\rho_2;q)_r}
\frac1{(q;q)_{n-r}(A;q)_{n+r}}.
\]
This is precisely $\alpha_r'/((q;q)_{n-r}(aq;q)_{n+r})$.  Summing over $r$ proves that $(\alpha',\beta')$ satisfies \eqref{eq:bailey-pair}.
\end{proof}

\begin{corollary}[The limiting Bailey-chain step]\label{cor:bailey-S}
If $(\alpha,\beta)$ is a Bailey pair relative to $a$, then so is $(\alpha^{S},\beta^{S})$, where
\begin{align}
\alpha_n^{S}&=a^nq^{n^2}\alpha_n, \label{eq:S-alpha}\\
\beta_n^{S}&=\sum_{j=0}^{n}\frac{a^jq^{j^2}}{(q;q)_{n-j}}\beta_j. \label{eq:S-beta}
\end{align}
\end{corollary}

\begin{proof}
Let $\rho_1,\rho_2\to\infty$ in Bailey's lemma.  Since
\[
(\rho;q)_n\sim(-\rho)^nq^{\binom n2},
\]
the factors in \eqref{eq:bailey-alpha-prime} and \eqref{eq:bailey-beta-prime} have exactly the limits displayed above.
\end{proof}

\begin{corollary}[Limiting Bailey identity]\label{cor:limiting-bailey}
For a Bailey pair relative to $a$, under absolute convergence,
\begin{equation}\label{eq:limiting-bailey}
\sum_{n=0}^{\infty}a^nq^{n^2}\beta_n
=\frac1{(aq;q)_\infty}
\sum_{n=0}^{\infty}a^nq^{n^2}\alpha_n.
\end{equation}
\end{corollary}

\begin{proof}
Apply \cref{cor:bailey-S}.  From the defining Bailey relation for $(\alpha^S,\beta^S)$,
\[
\beta_N^S=\sum_{r=0}^{N}
\frac{a^rq^{r^2}\alpha_r}{(q;q)_{N-r}(aq;q)_{N+r}}.
\]
As $N\to\infty$, the left side of \eqref{eq:S-beta} tends
\(
(q;q)_\infty^{-1}\sum_{j\ge0}a^jq^{j^2}\beta_j\), while the last display tends
\(
((q;q)_\infty(aq;q)_\infty)^{-1}
\sum_{r\ge0}a^rq^{r^2}\alpha_r
\).
Cancel $(q;q)_\infty^{-1}$.
\end{proof}

\subsection{A one-step Bailey lattice shift}

\begin{lemma}[Parameter-lowering shift]\label{lem:bailey-L}
If $(\alpha,\beta)$ is a Bailey pair relative to $a$, define $\alpha_{-1}=0$ and
\begin{equation}\label{eq:L-alpha}
\widetilde\alpha_n=(1-a)
\left(
\frac{\alpha_n}{1-aq^{2n}}
-\frac{aq^{2n-2}\alpha_{n-1}}{1-aq^{2n-2}}
\right),
\qquad
\widetilde\beta_n=\beta_n.
\end{equation}
Then $(\widetilde\alpha,\widetilde\beta)$ is a Bailey pair relative to $a/q$.
\end{lemma}

\begin{proof}
A pair relative to $a/q$ has denominator $(q;q)_{n-j}(a;q)_{n+j}$.  Substitute \eqref{eq:L-alpha} and reindex the second sum.  The coefficient of $\alpha_j$ is
\begin{align*}
&\frac{(1-a)(1-aq^{n+j})}
{(q;q)_{n-j}(a;q)_{n+j+1}(1-aq^{2j})}
-\frac{(1-a)aq^{2j}(1-q^{n-j})}
{(q;q)_{n-j}(a;q)_{n+j+1}(1-aq^{2j})}\\
&\qquad=\frac{1-a}{(q;q)_{n-j}(a;q)_{n+j+1}}
=\frac1{(q;q)_{n-j}(aq;q)_{n+j}}.
\end{align*}
The numerator in the first line simplifies because
\[
(1-aq^{n+j})-aq^{2j}(1-q^{n-j})=1-aq^{2j}.
\]
The resulting sum is exactly the original $\beta_n$.
\end{proof}

\begin{remark}
The operator in \cref{cor:bailey-S} moves along a Bailey chain while keeping $a$ fixed.  The shift in \cref{lem:bailey-L} changes $a$ to $a/q$.  Combining these two elementary operations is enough to prove every case of the Andrews--Gordon identities; no black-box multiple transformation will be used.
\end{remark}

\section{The Rogers--Ramanujan identities}

Define
\begin{equation}\label{eq:GH-def}
G(q)=\sum_{n=0}^{\infty}\frac{q^{n^2}}{(q;q)_n},
\qquad
H(q)=\sum_{n=0}^{\infty}\frac{q^{n^2+n}}{(q;q)_n}.
\end{equation}

\begin{theorem}[Rogers--Ramanujan identities]\label{thm:RR}
For $|q|<1$,
\begin{align}
G(q)&=\frac1{(q,q^4;q^5)_\infty}, \label{eq:RR1}\\
H(q)&=\frac1{(q^2,q^3;q^5)_\infty}. \label{eq:RR2}
\end{align}
\end{theorem}

\begin{proof}
Start from the unit Bailey pair and apply one Bailey-chain step \cref{cor:bailey-S}.  Since $\beta_n^{\mathrm U}=\delta_{n,0}$, the new beta sequence is
\begin{equation}\label{eq:unit-after-S-beta}
\beta_n^S=\frac1{(q;q)_n},
\end{equation}
while $\alpha_n^S=a^nq^{n^2}\alpha_n^{\mathrm U}(a)$.  Apply the limiting identity \eqref{eq:limiting-bailey} to this new pair.

For $a=1$, the removable limit in \eqref{eq:unit-alpha} is
\[
\alpha_0^{\mathrm U}(1)=1,
\qquad
\alpha_n^{\mathrm U}(1)=(-1)^nq^{\binom n2}(1+q^n)
\quad(n\ge1).
\]
Hence
\begin{align*}
G(q)
&=\frac1{(q;q)_\infty}
\left(1+\sum_{n\ge1}(-1)^n
q^{(5n^2-n)/2}(1+q^n)\right)\\
&=\frac1{(q;q)_\infty}
\sum_{n\in\IntegerNumbers}(-1)^nq^{(5n^2-n)/2}.
\end{align*}
The second equality pairs the positive index $n$ with the negative index $-n$.  Apply the Jacobi triple product with base $q^5$ and $z=q^2$:
\[
\sum_{n\in\IntegerNumbers}(-1)^nq^{(5n^2-n)/2}
=(q^5,q^2,q^3;q^5)_\infty.
\]
Dividing by
\[
(q;q)_\infty=(q,q^2,q^3,q^4,q^5;q^5)_\infty
\]
gives \eqref{eq:RR1}.

For $a=q$, the unit alpha is
\[
\alpha_n^{\mathrm U}(q)=(-1)^nq^{\binom n2}
\frac{1-q^{2n+1}}{1-q}.
\]
The same calculation gives
\begin{align*}
H(q)
&=\frac1{(q^2;q)_\infty}
\sum_{n\ge0}q^{2n^2+2n}\alpha_n^{\mathrm U}(q)\\
&=\frac1{(q;q)_\infty}
\sum_{n\in\IntegerNumbers}(-1)^nq^{(5n^2+3n)/2}.
\end{align*}
Here the negative index $-n-1$ supplies the term obtained by multiplying the $n$th nonnegative term by $-q^{2n+1}$.  The triple product with base $q^5$ and $z=q^4$ gives
\[
\sum_{n\in\IntegerNumbers}(-1)^nq^{(5n^2+3n)/2}
=(q^5,q^4,q;q^5)_\infty.
\]
Division by $(q;q)_\infty$ yields \eqref{eq:RR2}.
\end{proof}

\subsection{The partition theorem encoded by the two series}

\begin{corollary}[Rogers--Ramanujan partition theorem]\label{cor:RR-partitions}
For every $N\ge0$:
\begin{enumerate}[label=(\roman*)]
\item the number of partitions of $N$ whose adjacent parts differ by at least $2$ equals the number of partitions into parts congruent to $1$ or $4$ modulo $5$;
\item the number of partitions of $N$ whose adjacent parts differ by at least $2$ and whose smallest part is at least $2$ equals the number of partitions into parts congruent to $2$ or $3$ modulo $5$.
\end{enumerate}
\end{corollary}

\begin{proof}
A partition with exactly $n$ parts and successive differences at least $2$ contains the staircase
\[
2n-1,\ 2n-3,\ \ldots,\ 3,\ 1,
\]
of total size $n^2$.  Subtracting it leaves a partition with at most $n$ parts, whose generating function is $1/(q;q)_n$.  Summing over $n$ gives $G(q)$.  If the smallest part is at least $2$, subtract instead
\[
2n,\ 2n-2,\ \ldots,\ 4,\ 2,
\]
of total size $n^2+n$, giving $H(q)$.  The product sides of \eqref{eq:RR1} and \eqref{eq:RR2} are respectively the generating functions for unrestricted repetitions of parts in the indicated residue classes.  Equality of coefficients proves both statements.
\end{proof}

\section{The Rogers--Ramanujan continued fraction}

For $0<|q|<1$, fix a branch of $q^{1/5}$ and define
\begin{equation}\label{eq:R-def}
R(q)=q^{1/5}\frac{H(q)}{G(q)}.
\end{equation}
The product formulas immediately imply
\begin{equation}\label{eq:R-product}
R(q)=q^{1/5}\frac{(q,q^4;q^5)_\infty}{(q^2,q^3;q^5)_\infty}.
\end{equation}
We now prove that this quotient is the classical continued fraction.

\begin{theorem}[Rogers--Ramanujan continued fraction]\label{thm:RRCF}
For $|q|<1$,
\begin{equation}\label{eq:RRCF}
R(q)=\cfrac{q^{1/5}}{
1+\cfrac{q}{
1+\cfrac{q^2}{
1+\cfrac{q^3}{1+\ddots}}}}.
\end{equation}
\end{theorem}

\begin{proof}
Let
\[
C_N=1+\cfrac{q}{1+\cfrac{q^2}{\ddots+\cfrac{q^N}{1}}}
=\frac{P_N}{Q_N}
\]
be the $N$th finite convergent.  Continuant expansion gives
\begin{align}
P_{-1}=1,\qquad P_0=1,\qquad P_N=P_{N-1}+q^NP_{N-2}, \label{eq:P-recurrence}\\
Q_{-1}=0,\qquad Q_0=1,\qquad Q_N=Q_{N-1}+q^NQ_{N-2}. \label{eq:Q-recurrence}
\end{align}
We claim
\begin{align}
P_N&=\sum_{j=0}^{\lfloor(N+1)/2\rfloor}
q^{j^2}\GaussianBinomial{N+1-j}{j}{q}, \label{eq:P-polynomial}\\
Q_N&=\sum_{j=0}^{\lfloor N/2\rfloor}
q^{j^2+j}\GaussianBinomial{N-j}{j}{q}. \label{eq:Q-polynomial}
\end{align}
The initial cases are immediate.  Substitution of the $q$-Pascal recurrence \eqref{eq:qpascal1} verifies that the right sides satisfy \eqref{eq:P-recurrence} and \eqref{eq:Q-recurrence}; uniqueness of the recurrence proves the claim.

For fixed $j$,
\[
\GaussianBinomial{N+1-j}{j}{q}
=\frac{(q^{N+2-2j};q)_j}{(q;q)_j}\longrightarrow\frac1{(q;q)_j},
\]
and similarly for \eqref{eq:Q-polynomial}.  On every compact disk $|q|\le r<1$, these terms are bounded by a constant times $r^{j^2}/|(q;q)_j|$, so dominated convergence applies.  Therefore
\[
P_N\longrightarrow G(q),
\qquad
Q_N\longrightarrow H(q).
\]
The limits are nonzero near $q=0$ and hence, by the product formulas, throughout the disk away from their product zeros; in fact those products have no zeros for $|q|<1$.  Consequently
\[
C_N\longrightarrow\frac{G(q)}{H(q)}.
\]
Taking reciprocals and multiplying by $q^{1/5}$ gives \eqref{eq:RRCF}.
\end{proof}

\begin{remark}
The finite polynomials in \eqref{eq:P-polynomial}--\eqref{eq:Q-polynomial} are Schur-type finitizations.  They make the continued fraction, the Gaussian polynomial, and the Rogers--Ramanujan series three views of the same recurrence.
\end{remark}

\section{The full Andrews--Gordon identities}

The supplied Andrews--Gordon entry presents a multisum whose product modulus is $2k+1$.  The usual Bailey chain proves the two extreme residue choices immediately; the parameter-lowering shift of \cref{lem:bailey-L} inserts the linear tail required for every intermediate choice.

\begin{lemma}[Iterating the limiting Bailey step]\label{lem:iterated-S}
Start with $\beta_n=\delta_{n,0}$ and apply the transformation $S_a$ of \cref{cor:bailey-S} exactly $t\ge1$ times.  Then
\begin{equation}\label{eq:iterated-beta}
\beta_n^{(t)}=
\sum_{n\ge s_1\ge\cdots\ge s_{t-1}\ge0}
\frac{a^{s_1+\cdots+s_{t-1}}
 q^{s_1^2+\cdots+s_{t-1}^2}}
{(q;q)_{n-s_1}(q;q)_{s_1-s_2}\cdots
 (q;q)_{s_{t-2}-s_{t-1}}(q;q)_{s_{t-1}}},
\end{equation}
with the empty-chain interpretation $\beta_n^{(1)}=1/(q;q)_n$.  The corresponding alpha sequence is
\begin{equation}\label{eq:iterated-alpha}
\alpha_n^{(t)}=a^{tn}q^{tn^2}\alpha_n.
\end{equation}
\end{lemma}

\begin{proof}
Formula \eqref{eq:iterated-alpha} follows immediately by multiplying the alpha factor in \eqref{eq:S-alpha} at each step.  The case $t=1$ of \eqref{eq:iterated-beta} follows from $\beta_j=\delta_{j,0}$.  If it holds for $t$, substituting it in \eqref{eq:S-beta} adds a new outer index and a new denominator $(q;q)_{n-s_1}$, giving the formula for $t+1$.
\end{proof}

\begin{theorem}[Andrews--Gordon identity]\label{thm:AG}
Let $k\ge2$ and $1\le i\le k$.  For nonnegative integers $n_1,\ldots,n_{k-1}$ put
\[
N_j=n_j+n_{j+1}+\cdots+n_{k-1}\qquad(1\le j\le k-1).
\]
Then, for $|q|<1$,
\begin{equation}\label{eq:AG}
\sum_{n_1,\ldots,n_{k-1}\ge0}
\frac{q^{N_1^2+\cdots+N_{k-1}^2+N_i+\cdots+N_{k-1}}}
{(q;q)_{n_1}\cdots(q;q)_{n_{k-1}}}
=
\frac{(q^i,q^{2k+1-i},q^{2k+1};q^{2k+1})_\infty}
{(q;q)_\infty},
\end{equation}
where the linear sum is empty when $i=k$.
\end{theorem}

\begin{proof}
Set
\[
K=k-1,
\qquad M=2K+3=2k+1.
\]
It is convenient to use decreasing variables
\[
s_j=N_j\quad(1\le j\le K),
\qquad s_{K+1}=0.
\]
Then $n_j=s_j-s_{j+1}$, and the left side of \eqref{eq:AG} is
\begin{equation}\label{eq:AG-monotone}
\sum_{s_1\ge\cdots\ge s_K\ge0}
\frac{q^{s_1^2+\cdots+s_K^2+s_i+\cdots+s_K}}
{(q;q)_{s_1-s_2}\cdots(q;q)_{s_{K-1}-s_K}(q;q)_{s_K}}.
\end{equation}
We derive this sum from the unit Bailey pair relative to $q$.

\medskip
\noindent\textbf{Case 1: $i=1$.}
Apply $S_q$ exactly $K$ times, and then apply the limiting identity \eqref{eq:limiting-bailey}, still relative to $q$.  By \cref{lem:iterated-S}, the left side is exactly \eqref{eq:AG-monotone}: the final limiting sum supplies the outer variable $s_1$ and its weight $q^{s_1^2+s_1}$, while the preceding $K-1$ nested variables have the same quadratic-plus-linear weights.

On the alpha side, the $K$ chain steps and the final limiting identity together multiply the unit alpha by $q^{(K+1)(n^2+n)}$.  Using \eqref{eq:unit-alpha},
\begin{align*}
\text{RHS}
&=\frac1{(q^2;q)_\infty}
\sum_{n\ge0}q^{(K+1)(n^2+n)}
(-1)^nq^{\binom n2}\frac{1-q^{2n+1}}{1-q}\\
&=\frac1{(q;q)_\infty}
\sum_{m\in\IntegerNumbers}(-1)^m
q^{\{Mm^2+(M-2)m\}/2}.
\end{align*}
The equality in the second line pairs the indices $m=n\ge0$ and $m=-n-1$; the $n=0$ pair contributes $1-q$.  The Jacobi triple product with base $q^M$ and $z=q^{M-1}$ evaluates the bilateral sum as
\[
(q^M,q^{M-1},q;q^M)_\infty.
\]
This is the product in \eqref{eq:AG} for $i=1$.

\medskip
\noindent\textbf{Case 2: $2\le i\le K+1=k$.}
Put
\[
t=K-i+2,
\qquad u=i-2,
\]
so that $t\ge1$, $u\ge0$, and $t+u=K$.  Starting with the unit pair relative to $q$, perform the following operations:
\[
\underbrace{S_q\circ\cdots\circ S_q}_{t\text{ times}},
\qquad L_q,
\qquad
\underbrace{S_1\circ\cdots\circ S_1}_{u\text{ times}},
\]
where $L_q$ is the shift in \cref{lem:bailey-L}; finally apply the limiting identity relative to $1$.

The first $S_q$ applied to the delta beta sequence creates only the terminal factor $1/(q;q)_n$.  The remaining $t-1$ applications create the innermost variables, each weighted by $q^{s_j^2+s_j}$.  The shift $L_q$ leaves beta unchanged and changes the relative parameter to $1$.  The next $u$ applications create variables weighted by $q^{s_j^2}$, and the final limiting identity creates one more outer variable with weight $q^{s_1^2}$.  Since $u=i-2$, the variables with linear weight begin exactly at position $i$.  Thus the left side is \eqref{eq:AG-monotone}.

It remains to calculate the alpha side.  After the first $t$ chain steps, put
\[
\gamma_n=q^{t(n^2+n)}\alpha_n^{\mathrm U}(q).
\]
Applying $L_q$ gives $\delta_0=1$ and, for $n\ge1$,
\begin{align}
\delta_n
&=(1-q)\left(
\frac{\gamma_n}{1-q^{2n+1}}
-\frac{q^{2n-1}\gamma_{n-1}}{1-q^{2n-1}}
\right) \notag\\
&=(-1)^nq^{\{(2t+1)n^2-(2t-1)n\}/2}
\left(1+q^{(2t-1)n}\right). \label{eq:delta-AG}
\end{align}
This is a direct substitution of the unit alpha formula; the two exponents differ by $(2t-1)n$.

The subsequent $u$ applications of $S_1$ and the final limiting identity multiply \eqref{eq:delta-AG} by $q^{(u+1)n^2}=q^{(i-1)n^2}$.  Since
\[
2(t+i)-1=2K+3=M,
\qquad
2t-1=M-2i,
\]
the alpha side becomes
\begin{align*}
\frac1{(q;q)_\infty}
\left(
1+\sum_{n\ge1}(-1)^n
q^{\{Mn^2-(M-2i)n\}/2}
(1+q^{(M-2i)n})
\right)
&=\frac1{(q;q)_\infty}
\sum_{m\in\IntegerNumbers}(-1)^m
q^{\{Mm^2-(M-2i)m\}/2}\\
&=\frac{(q^M,q^i,q^{M-i};q^M)_\infty}
{(q;q)_\infty},
\end{align*}
where the first equality pairs $m=n$ and $m=-n$, and the second is the Jacobi triple product with $z=q^i$.  This is precisely \eqref{eq:AG}.
\end{proof}

\begin{corollary}
When $k=2$, the cases $i=2$ and $i=1$ of \cref{thm:AG} reduce respectively to the first and second Rogers--Ramanujan identities.
\end{corollary}

\begin{proof}
There is one variable $N_1=n_1$.  For $i=2$ the linear term is absent and the sum is $G(q)$; for $i=1$ it is $H(q)$.  The two products have the residue pairs $\{1,4\}$ and $\{2,3\}$ modulo $5$ after cancellation against $(q;q)_\infty$.
\end{proof}

\section{How the remaining supplied identities fit}

This section is an atlas, not a second list of theorems.  Its purpose is to locate every supplied file in the proof architecture developed above.  Formulas mentioned here are not used subsequently; the fully proved results are those stated as theorems in the preceding sections.

\begin{longtable}{@{}p{0.25\textwidth}p{0.28\textwidth}p{0.40\textwidth}@{}}
\toprule
\textbf{Supplied topic} & \textbf{Primary mechanism} & \textbf{Relation to this article}\\
\midrule
\endfirsthead
\toprule
\textbf{Supplied topic} & \textbf{Primary mechanism} & \textbf{Relation to this article}\\
\midrule
\endhead
$q$-Pochhammer symbol; $q$-binomial coefficient
& Shifted-factorial algebra; Gaussian recurrences
& Definitions and complete finite/infinite $q$-binomial proofs are in \cref{prop:q-pascal,thm:finite-qbinomial,thm:infinite-qbinomial}.\\[0.4em]
$q$-hypergeometric function; $q$-Vandermonde; $q$-Saalsch\"utz
& Heine transformation; balanced terminating summation
& The modern ${}_r\phi_s$ normalization, Heine transformation, $q$-Gauss, $q$-Chu--Vandermonde, and $q$-Pfaff--Saalsch\"utz are proved in \cref{thm:heine,cor:q-gauss,cor:q-chu,thm:q-pfaff}.\\[0.4em]
Fine's equation
& Rearrangement of a ${}_2\phi_1$ or a specialization of Heine transformations
& It belongs to the same sum-interchange calculus as \cref{thm:heine}.  Fine's identities are especially effective when one parameter is sent to $0$ or infinity.\\[0.4em]
Jackson's identity; Jackson--Slater identity
& Terminating very-well-poised basic hypergeometric transformations
& These lie above $q$-Pfaff--Saalsch\"utz in the hierarchy.  Their inner balanced sums are evaluated by the same telescoping or Bailey-transform principles.\\[0.4em]
Dougall--Ramanujan identity
& Very-well-poised ${}_8\phi_7$ summation and limiting transformations
& It is a high-level master summation: many lower identities arise by specialization.  Bhatnagar gives an Euler-telescoping proof of Jackson's $q$-Dougall sum \cite{Bhatnagar2011}; the present article proves the lower balanced theorem needed for Bailey's lemma.\\[0.4em]
Jacobi triple product
& Laurent functional equation plus a constant-term evaluation
& Proved completely in \cref{thm:jtp}; this proof combines Euler's expansions with the Durfee-square identity.\\[0.4em]
Quintuple product identity
& Combination and residue dissection of triple products
& A five-factor product is converted into a difference of two theta progressions.  The proof engine is the same functional equation and residue splitting used in \cref{thm:jtp,cor:pentagonal}.\\[0.4em]
Schr\"oter's formula
& Bilinear addition law for Ramanujan theta functions
& Ramanujan's general theta product is proved in \cref{cor:ramanujan-theta}.  Schr\"oter-type formulas follow by multiplying two bilateral theta series and making an invertible linear change of the two summation indices.\\[0.4em]
Jacobi theta functions; Ramanujan theta functions
& Triple products; Poisson summation
& Product formulas and modular inversion are proved in \cref{cor:ramanujan-theta,thm:theta-products,thm:theta-transform}.\\[0.4em]
Dedekind eta function
& Theta product identity and modular inversion
& The two generators $\tau\mapsto\tau+1$ and $\tau\mapsto-1/\tau$ are proved in \cref{thm:eta-transform}.\\[0.4em]
Partition functions $P$ and $Q$; pentagonal theorem
& Euler products; Ferrers diagrams; triple-product specialization
& The ordinary product, distinct-part product, odd/distinct equivalence, Durfee decomposition, pentagonal theorem, and recurrence are proved in \cref{thm:partition-product,thm:distinct-product,thm:durfee,cor:pentagonal,cor:partition-recurrence}.\\[0.4em]
Sum of squares function
& Powers of theta series and modular/Lambert expansions
& The universal generating function is proved in \cref{prop:squares-theta}.  Explicit divisor-sum formulas arise from further modular-form decompositions.\\[0.4em]
Rogers--Ramanujan identities
& Unit Bailey pair; one chain step; triple product
& Both analytic identities and both partition interpretations are proved in \cref{thm:RR,cor:RR-partitions}.\\[0.4em]
Rogers--Ramanujan continued fraction
& Continuant recurrence; Schur polynomial finitization
& The product and continued-fraction equality are proved in \cref{thm:RRCF}; the finite convergents use Gaussian polynomials.\\[0.4em]
Andrews--Gordon identity
& Bailey chain plus one parameter-lowering step
& Every $1\le i\le k$ case is proved in \cref{thm:AG}; the intermediate cases are obtained without assuming a black-box multiple transformation.\\[0.4em]
Rogers--Selberg; Bailey mod $9$; Rogers mod $14$; Dyson mod $27$
& Specialized Bailey pairs, Bailey chains/lattices, and final theta-product recognition
& These are higher-modulus members of the same workflow as \cref{thm:RR,thm:AG}.  The modulus is read from the quadratic exponent in the final bilateral theta series.\\[0.4em]
General $q$-series identities notebook
& Product dissections, Lambert series, modular equations
& Its identities combine the article's basic transformations with residue-class extraction.  The triple product explains why apparently unrelated sums collapse to short products.\\[0.4em]
Borwein conjectures
& Coefficient sign patterns in finite products; saddle-point asymptotics and positivity-preserving transforms
& The historical MathWorld page predates major progress.  The first conjecture was proved by Wang \cite{Wang2022}; Krattenthaler and Wang proved the first and second and obtained a partial cubic result \cite{KrattenthalerWang2022}.  New modulus-$3$ and modulus-$5$ conjectures continue the program \cite{BerkovichDhar2025}.\\
\bottomrule
\end{longtable}

\subsection{A reusable discovery protocol}

The atlas suggests a practical protocol for proving or discovering identities of the type found in the archive:
\begin{enumerate}[label=\textbf{Step \arabic*.}]
\item Normalize every finite product with $q$-Pochhammer notation and identify whether the series is terminating, balanced, well-poised, or very-well-poised.
\item Search first for a $q$-difference equation, a finite telescoping certificate, or a Heine transformation.  These are less expensive than guessing a product.
\item If a sum has a denominator $(q;q)_n$ and a quadratic exponent, test whether it is the beta side of a Bailey pair.  Use Bailey inversion to recover a candidate alpha sequence.
\item Convert the final alpha sum to a bilateral theta series.  The coefficient of $n^2$ predicts the product modulus, while the coefficient of $n$ predicts the excluded residue classes.
\item For finite positivity problems, keep the Gaussian-polynomial form instead of taking $N\to\infty$ too early.  Positivity, reciprocity, and unimodality are finite phenomena and may disappear in the limiting product notation.
\end{enumerate}

\section{Further deductions and research directions}

The proofs above support several concrete extensions.

\subsection{Finite forms before infinite limits}

The continued-fraction proof shows the value of retaining the finite Gaussian-polynomial approximants.  The same principle applies to the Rogers--Ramanujan and Andrews--Gordon sums: before taking the limiting Bailey identity, stop at a finite outer index.  One obtains polynomial identities whose limits are the infinite products.  Such finitizations are particularly promising for coefficientwise inequalities, truncated pentagonal identities, and computer-assisted discovery, because every claim becomes a finite statement in $\IntegerNumbers[q]$.

\subsection{Automating Bailey-pair discovery}

Bailey inversion \eqref{eq:bailey-inversion} turns pair discovery into a deterministic calculation.  Given a plausible beta sequence, compute its alpha sequence, factor the result, and inspect whether it has a short theta-like form.  If it does, iterating $S_a$ or inserting the shift $L_a$ will produce families.  The proof remains human-readable because the only verification is the already proved triangular inversion.

\subsection{Modulus prediction from quadratic forms}

In the Rogers--Ramanujan and Andrews--Gordon proofs, the final bilateral sum has the form
\[
\sum_{n\in\IntegerNumbers}(-1)^nq^{(Mn^2+Bn)/2}.
\]
The Jacobi product forces modulus $M$, and the linear coefficient determines the two nontrivial factors.  This gives a useful reverse-engineering rule: a multisum whose Bailey reduction produces a quadratic coefficient $M/2$ should be tested against products modulo $M$.  The mod $9$, $14$, and $27$ files in the archive are natural test cases for systematic implementations of this rule.

\subsection{Borwein-type finite positivity}

The modern proofs of the first and second Borwein conjectures are asymptotic and analytic, while many extensions remain conjectural.  The finite $q$-binomial machinery in \cref{thm:finite-qbinomial,lem:q-annihilation,thm:bailey-inversion} suggests a complementary program: search for sign-preserving triangular transforms that map a difficult finite dissection to manifestly nonnegative Gaussian-polynomial sums.  The key constraint is to preserve the residue-class decomposition rather than merely the total polynomial.

\subsection{Bilateral and multilateral Bailey structures}

Recent work has clarified bilateral Bailey lattices and new Andrews--Gordon-type identities \cite{DousseJouhetKonan2025,DousseJangJouhet2025}.  The elementary proof of \cref{lem:bailey-L} makes clear why bilateral extensions are plausible: the parameter shift is a first-order difference of neighboring alpha terms.  Allowing indices in $\IntegerNumbers$ replaces the boundary condition $\alpha_{-1}=0$ by a convergence condition at $-\infty$, opening a route to mock theta, Appell--Lerch, and false-theta identities.

\section{Conclusion}

A surprisingly small toolkit accounts for the large collection supplied with this project.  The finite $q$-binomial theorem controls Gaussian polynomials and finite products.  Its infinite form proves Euler expansions and Heine transformations.  A single telescoping lemma proves the balanced $q$-Saalsch\"utz sum.  Ferrers diagrams evaluate the constant term needed for the Jacobi triple product.  Poisson summation supplies modular inversion.  Finally, Bailey inversion and two elementary pair transformations turn isolated Rogers--Ramanujan phenomena into the full Andrews--Gordon family.

The main conceptual lesson is that the subject alternates between two languages.  Products make congruence restrictions and modular structure visible; sums make recurrences, partitions, and transformations visible.  The decisive proofs are the dictionaries between those languages.

\appendix

\section{Convergence and interchange ledger}

For reference, we record the analytic justifications used in the main text.

\begin{enumerate}[label=\textbf{A.\arabic*.}]
\item \textbf{Infinite products.}  \Cref{lem:product-convergence} gives uniform convergence on compact parameter sets.  Ratios of products are therefore holomorphic away from explicit zeros of denominator factors.
\item \textbf{Euler and basic hypergeometric series.}  For $|z|<1$, the coefficient ratio in the infinite $q$-binomial series tends to $z$.  In Heine's proof, one may initially impose $|z|<1$, $|b|<1$, and compact parameter bounds; the resulting double series is absolutely convergent.  The final identity extends by analytic continuation.
\item \textbf{Jacobi's bilateral series.}  On a compact annulus $r\le|z|\le R$, the positive and negative tails are bounded by geometric multiples of $|q|^{n^2/3}$ for large $|n|$.  Hence normal convergence holds on $\ComplexNumbers^\times$.
\item \textbf{Partition products.}  Every coefficient of $q^N$ uses only factors indexed by $m\le N$, so all combinatorial product manipulations are valid in $\IntegerNumbers[[q]]$ independently of analytic convergence.
\item \textbf{Poisson summation.}  Schwartz decay permits uniform differentiation of the periodization, termwise integration, and absolute convergence of the Fourier series.
\item \textbf{Bailey limits.}  The Rogers--Ramanujan and Andrews--Gordon applications have summands bounded by $C|q|^{cn^2}$ times products whose reciprocals remain bounded for $|q|\le r<1$.  Thus the limiting operations and sum rearrangements are absolutely and locally uniformly convergent.
\item \textbf{Continued fractions.}  The explicit convergent polynomials are dominated by the convergent Rogers--Ramanujan series.  Their numerator and denominator limits are nonzero products in the unit disk, so passage to the quotient is valid.
\end{enumerate}

\section{Formula sheet}

For convenient reuse, the central identities are collected here without proof references suppressed:
\begin{align*}
(a;q)_{m+n}&=(a;q)_m(aq^m;q)_n,\\
(z;q)_N&=\sum_{k=0}^{N}(-1)^kq^{\binom k2}\GaussianBinomial{N}{k}{q} z^k,\\
\sum_{n\ge0}\frac{(a;q)_n}{(q;q)_n}z^n
&=\frac{(az;q)_\infty}{(z;q)_\infty},\\
\sum_{n\in\IntegerNumbers}(-1)^nq^{\binom n2}z^n
&=(q,z,q/z;q)_\infty,\\
\sum_{n\ge0}\frac{q^{n^2}}{(q;q)_n}
&=\frac1{(q,q^4;q^5)_\infty},\\
\sum_{n\ge0}\frac{q^{n^2+n}}{(q;q)_n}
&=\frac1{(q^2,q^3;q^5)_\infty},\\
\sum_{n_1,\ldots,n_{k-1}\ge0}
\frac{q^{\sum_{j=1}^{k-1}N_j^2+\sum_{j=i}^{k-1}N_j}}
{\prod_{j=1}^{k-1}(q;q)_{n_j}}
&=\frac{(q^i,q^{2k+1-i},q^{2k+1};q^{2k+1})_\infty}
{(q;q)_\infty}.
\end{align*}

\section{Proof-completeness index}

\begin{longtable}{@{}p{0.43\textwidth}p{0.48\textwidth}@{}}
\toprule
\textbf{Result} & \textbf{Complete proof location}\\
\midrule
\endfirsthead
\toprule
\textbf{Result} & \textbf{Complete proof location}\\
\midrule
\endhead
Finite $q$-binomial and reciprocal expansions & \Cref{thm:finite-qbinomial,cor:finite-reciprocal}\\
Infinite $q$-binomial and Euler expansions & \Cref{thm:infinite-qbinomial,cor:euler-expansions}\\
Heine, $q$-Gauss, and $q$-Chu--Vandermonde & \Cref{thm:heine,cor:q-gauss,cor:q-chu}\\
$q$-Pfaff--Saalsch\"utz & \Cref{lem:euler-telescope,thm:q-pfaff}\\
Partition product, distinct/odd, and Durfee square & \Cref{thm:partition-product,thm:distinct-product,thm:durfee}\\
Jacobi triple product and Euler pentagonal theorem & \Cref{thm:jtp,cor:pentagonal}\\
Theta products, Poisson summation, theta and eta modular laws & \Cref{thm:theta-products,thm:poisson,thm:theta-transform,thm:eta-transform}\\
Bailey inversion, unit pair, Bailey lemma, chain and lattice shift & \Cref{thm:bailey-inversion,cor:unit-pair,thm:bailey-lemma,cor:bailey-S,lem:bailey-L}\\
Rogers--Ramanujan identities and partition theorem & \Cref{thm:RR,cor:RR-partitions}\\
Rogers--Ramanujan continued fraction & \Cref{thm:RRCF}\\
All Andrews--Gordon cases & \Cref{lem:iterated-S,thm:AG}\\
\bottomrule
\end{longtable}

\begin{thebibliography}{99}

\bibitem{AAB1987}
A.~K. Agarwal, G.~E. Andrews, and D.~M. Bressoud,
\newblock The Bailey lattice,
\newblock \emph{J. Indian Math. Soc.} \textbf{51} (1987), 57--73.

\bibitem{Andrews1974}
G.~E. Andrews,
\newblock An analytic generalization of the Rogers--Ramanujan identities for odd moduli,
\newblock \emph{Proc. Natl. Acad. Sci. USA} \textbf{71} (1974), 4082--4085.

\bibitem{Andrews1976}
G.~E. Andrews,
\newblock \emph{The Theory of Partitions},
\newblock Addison--Wesley, 1976; Cambridge University Press reprint, 1998.

\bibitem{Andrews1986}
G.~E. Andrews,
\newblock \emph{$q$-Series: Their Development and Application in Analysis, Number Theory, Combinatorics, Physics, and Computer Algebra},
\newblock CBMS Regional Conference Series 66, American Mathematical Society, 1986.

\bibitem{Bailey1949}
W.~N. Bailey,
\newblock Identities of the Rogers--Ramanujan type,
\newblock \emph{Proc. London Math. Soc.} (2) \textbf{50} (1949), 1--10.

\bibitem{BerkovichDhar2025}
A. Berkovich and A. Dhar,
\newblock New Borwein-type conjectures,
\newblock \emph{Experimental Mathematics} (2025); arXiv:2407.13788.

\bibitem{Bhatnagar2011}
G. Bhatnagar,
\newblock In praise of an elementary identity of Euler,
\newblock \emph{Electron. J. Combin.} \textbf{18}(2) (2011), Paper P13.

\bibitem{DLMF17}
NIST Digital Library of Mathematical Functions,
\newblock Chapter 17: $q$-hypergeometric and related functions,
\newblock \url{https://dlmf.nist.gov/17}, accessed 30 August 2026.

\bibitem{DLMF20}
NIST Digital Library of Mathematical Functions,
\newblock Chapter 20: Theta functions,
\newblock \url{https://dlmf.nist.gov/20}, accessed 30 August 2026.

\bibitem{DousseJangJouhet2025}
J. Dousse, J. Jang, and F. Jouhet,
\newblock Andrews--Gordon and Stanton type identities: bijective and Bailey lemma approaches,
\newblock arXiv:2507.13239 (2025).

\bibitem{DousseJouhetKonan2025}
J. Dousse, F. Jouhet, and I. Konan,
\newblock Bilateral Bailey lattices and Andrews--Gordon type identities,
\newblock \emph{SIGMA} \textbf{21} (2025), Paper 032.

\bibitem{GasperRahman2004}
G. Gasper and M. Rahman,
\newblock \emph{Basic Hypergeometric Series}, 2nd ed.,
\newblock Encyclopedia of Mathematics and its Applications 96, Cambridge University Press, 2004.

\bibitem{Gordon1961}
B. Gordon,
\newblock A combinatorial generalization of the Rogers--Ramanujan identities,
\newblock \emph{Amer. J. Math.} \textbf{83} (1961), 393--399.

\bibitem{GutierrezEtAl2025}
A. Guti\'errez, A.~L. Mart\'inez, M. Szwej, and M. Wildon,
\newblock A new bijective proof of the $q$-Pfaff--Saalsch\"utz identity with applications to quantum groups,
\newblock arXiv:2505.08422 (2025).

\bibitem{Jacobi1829}
C.~G.~J. Jacobi,
\newblock \emph{Fundamenta Nova Theoriae Functionum Ellipticarum},
\newblock K\"onigsberg, 1829.

\bibitem{KrattenthalerWang2022}
C. Krattenthaler and C. Wang,
\newblock An asymptotic approach to Borwein-type sign pattern theorems,
\newblock arXiv:2201.12415 (2022).

\bibitem{Rogers1894}
L.~J. Rogers,
\newblock Second memoir on the expansion of certain infinite products,
\newblock \emph{Proc. London Math. Soc.} \textbf{25} (1894), 318--343.

\bibitem{Wang2022}
C. Wang,
\newblock Analytic proof of the Borwein conjecture,
\newblock \emph{Adv. Math.} \textbf{394} (2022), Article 108028.

\bibitem{Warnaar2001}
S.~O. Warnaar,
\newblock 50 years of Bailey's lemma,
\newblock in \emph{Algebraic Combinatorics and Applications}, Springer, 2001, pp. 333--347.

\bibitem{MathWorldCollection}
E.~W. Weisstein and contributors,
\newblock MathWorld entries and associated Wolfram notebooks on $q$-series, theta functions, partitions, Rogers--Ramanujan identities, and related topics,
\newblock supplied in the project archive; current pages at \url{https://mathworld.wolfram.com/}, accessed 30 August 2026.

\end{thebibliography}

\end{document}
